% migrate-to: price-impact-execution
\subsection{Market Impact Models}

\begin{taxonomy} \hlt{Model Family Segmentation}
\begin{enumerate}[label=\roman*.]
\setlength{\itemsep}{0pt}
\item \hlt{Almgren-Chriss}: optimise a schedule by balancing expected impact cost against execution risk;
\item \hlt{linear impact}: use tractable instantaneous or temporary cost proportional to trading rate or signed size;
\item \hlt{propagator}: model each signed trade as a decaying impulse in future prices;
\item \hlt{empirical square-root}: summarise metaorder cost as concave in order size relative to market volume.
\end{enumerate}
\end{taxonomy}

\begin{principle} \hlt{Permanent Impact or Price Anticipation}\\
Long-horizon impact mixes metaorder pressure and information revealed by trading. Remove the information effect before calling the residual permanent impact.
\end{principle}

\begin{taxonomy} \hlt{Permanent-Impact View Segmentation}
\begin{enumerate}[label=\roman*.]
\setlength{\itemsep}{0pt}
\item \hlt{mechanical view}: prices move because buying and selling pressure are imbalanced;
\item \hlt{informational view}: prices move because the trade reveals new expectations about value;
\item \hlt{hybrid view}: executions contain both pressure and information, so identification matters more than labels.
\end{enumerate}
\end{taxonomy}

\begin{model} \hlt{Propagator Midprice Dynamics}\\
In trade time, a standard propagator writes the midprice as $m_k=m_{-\infty}+\sum\limits_{k'<k}G(k-k')\epsilon_{k'}+\eta_k$.\\
Here $m_k$ is the midprice after trade $k$, $\epsilon_k\in\{-1,+1\}$ is trade sign, $G(\cdot)$ is the decaying impact kernel, $\eta_k$ is fair-price noise, and $m_{-\infty}$ is an initial price level.
\end{model}

\begin{mechanism} \hlt{Trade-Sign Memory and Kernel Decay}\\
Trade signs are positively autocorrelated because large orders are split and similar agents trade together.\\
A decaying kernel offsets this persistence: impact must decay enough that predictable order flow does not imply predictable price drift.
\end{mechanism}

\begin{model} \hlt{Price Response Curve}\\
For horizon $h$, define the response curve by $\mathcal{R}(h)=\mathbb{E}[(m_{k+h}-m_k)\epsilon_k]$.\\
Here $\mathcal{R}(h)$ measures the expected signed price move after a trade. It is estimated with the trade-sign autocorrelation $C(\ell)=\mathbb{E}[\epsilon_{k+\ell}\epsilon_k]$ and the kernel increments $\Delta G(\ell)=G(\ell+1)-G(\ell)$.
\end{model}

\begin{taxonomy} \hlt{Propagator Kernel Segmentation}
\begin{enumerate}[label=\roman*.]
\setlength{\itemsep}{0pt}
\item \hlt{non-parametric}: estimate discrete kernel increments directly from response and sign autocorrelation;
\item \hlt{exponential}: impose fast decay with a small number of parameters;
\item \hlt{power law}: allow slow decay and long-memory order flow;
\item \hlt{agent-specific}: estimate separate kernels or intensities by participant class.
\end{enumerate}
\end{taxonomy}

\begin{process} \hlt{Debiasing Long-Horizon Impact}
\begin{enumerate}[label=\roman*.]
\setlength{\itemsep}{0pt}
\item classify trades by whether they are likely information-driven or liquidity-driven;
\item remove market and sector moves with a benchmark or factor model;
\item adjust for autocorrelated metaorder flow after the execution day;
\item measure the remaining post-execution drift over several days.
\end{enumerate}
\end{process}

\begin{model} \hlt{CAPM-Like Post-Execution Decomposition}\\
For event day $D$, a stock's post-execution move can be split into systematic and idiosyncratic parts:
\begin{equation*}
\log(P_d)-\log(P_{d-1})=
\beta\left[\log(I_d)-\log(I_{d-1})\right]+\Delta W_d.
\end{equation*}
Here $P_d$ is the stock close, $I_d$ the reference index, $\beta$ the index exposure, and $\Delta W_d$ the residual stock-specific move. Cumulative $\Delta W_d$ isolates the profile after removing index co-movement. The intuition is attribution: do not charge the execution for a market move that the whole index experienced.
\end{model}

\begin{caveat} \hlt{Autocorrelated Flow Bias}\\
Metaorder flow can remain signed and autocorrelated for days. If later same-direction executions are ignored, their impact may be falsely attributed to the original order.
\end{caveat}
